\section{The one-dimensional calculus}
\label{calc:section-calculus}

Throughout this Section, $\mathcal C$ is a scheme satisfying the working
package of Definition~\ref{lev:working-package}.  Its fixed presentation is
denoted by $F$, and $b$ is an infinite cardinal bounding the set of labels
and the finite amount of structural data in the scheme.  We use the
occurrence sorts of that Definition: a tuple in several vertices is a
tuple of tagged occurrences, and a map of systems commutes with every
specified structural map.  No choice of a hull or a basis is part of a
diagram.  A level is an eligible strong limit cardinal of countable
cofinality.  All further closure requirements below are requirements on
functions of bounded arities and fixed scheme data.
For the supply of unbounded refinements and larger auxiliary levels,
we use the eligible-subclass hypothesis of
Convention~\ref{lev:eligible-subclass}: the input class contains a fixed
definable unbounded subclass closed under suprema of strictly increasing
countable sequences.  This subclass is used only to supply levels;
the bounds below depend on fixed scheme data and bounded arities.

Stationary independence and local character are central to the theory of
good frames; see \cite[II.2.1]{Shelah2009}.  The assertions below are
proved from the working package for bounded tuples and finite systems.
Their role is to supply the independence and descent properties needed
for the next isolation construction.

A \emph{query} is an enumerated parameter set $Z$ together with the
diagram of the displayed tuple and that enumeration.  Repetitions in an
enumeration are allowed.  Saying that an answer is correct on $Z$ means
equality of the whole joint diagram, rather than equality only on finite
subtuples.  At a level $\mu$, every tuple or query called bounded has
cardinality less than $\mu$.

\subsection{Testing, placement, and amalgamation}

\begin{lemma}[Parameter orbits for a scheme]
\label{calc:parameter-orbits}
Fix an infinite $\sigma$.  If
$\chi>b+\sigma$ and $\chi^\sigma=\chi$, then, at every sufficiently closed
level greater than $\chi$, there are at most $\chi$ diagrams of
$\sigma$-tuples over a fixed parameter tuple of length at most $\chi$.
Moreover, natural diagrams are unchanged by simultaneous disjoint block
transpositions: within each of a family of disjoint convex regions,
interchange an adjoining block belonging exclusively to the first tuple
family with one belonging exclusively to the second, retaining each
block's internal order.
\end{lemma}

\begin{proof}
Present the parameter object on a well-order.  After choosing one term
representation of each parameter, its supporting positions form a
well-ordered set of size at most $\chi$.  There are at most $\chi$ slots
and positions relative to this set.  Labels, comparisons and coincidences
for at most $\sigma$ additional positions therefore have at most
$\chi^\sigma=\chi$ possibilities.  Equal profiles give equal natural
diagrams.  The diagram-isomorphism clause of the working package turns
this into the asserted orbit bound over the actual parameters.

Here are the hypotheses needed for the block argument in
Lemma~\ref{lev:block-transpositions}.  We have just proved its parameter
orbit bound.  Repeating the two blocks along an arbitrary order is
permitted after padding: the scheme has the union and common-upper-bound
clauses, so Lemma~\ref{lev:directed-unions} supplies the padded
presentations and their order arrows.  Naturalness compares the two
patterns at a common higher eligible level.  If the transposition changed
a diagram, the repeated patterns would distinguish cuts by whether the
first block preceded or followed a parameter block.  A binary tree of
least height $\epsilon$ with $2^\epsilon>\chi$ has at most $\chi$
nonterminal nodes and more than $\chi$ branches.  Divider points at those
nodes give at most $\chi$ parameter blocks and more than $\chi$ different
diagrams over them.  The preceding orbit bound gives a contradiction.
The repetition uses only a set of positions, so a higher eligible level
can contain the entire test.  This proves the assertion also when there
are boundedly many simultaneous transpositions.
\end{proof}

\begin{definition}[Natural parameter layouts]
\label{calc:natural-layout}
A natural parameter layout consists of finitely many well-ordered
segments, a fixed finite set of colors and their prescribed order, and
specified families of terms on specified collections of those positions.
Its parameters are all the designated terms; vertex sorts may be among
the designations.  The data must have the following monotonicity
property: after freezing bounded extra supports, increasing the ordinal
lengths of the remaining slots gives an embedding of layouts that
preserves the designated parameter terms.  Bounded frozen parameters
are permitted and count towards the input arity.  The layout, labels and
color conventions are fixed set data.

The whole object $F(H)$, with $H$ well-ordered, is a natural layout.
Unions of designated vertex sorts generated on specified segments are
further examples.  An arbitrary subset of the parameters in an unrelated
target chart is not asserted to be a natural layout.
\end{definition}

\begin{proposition}[Uniform testing for natural layouts]
\label{calc:testing}
\label{calc:layout-testing}
For each fixed natural layout there is a monotone cardinal function
$g(\sigma)$ with the following property.  Two placed comparison types of
tuples of length at most $\sigma$ over that layout are equal if and only
if they agree on parameter queries of size at most $g(\sigma)$.
The bound is independent of the eligible level.  The extra tuple supports
may lie in a non-well-ordered extension of the layout.  The assertion is
uniform when the fixed layout and frozen bounded data range over a fixed
set of descriptors of bounded size.
\end{proposition}

\begin{proof}
Amalgamate the orders containing the two placements over the parameter
order.  Freeze the labels of the two tuples, their combined ordered
support $T$, the coincidences with old positions, and the colors.  There
are set-many such frozen patterns at the prescribed bound.  The
unfrozen well-ordered parameter positions are described by a vector of
ordinal lengths, one coordinate for each strict slot and each prescribed
segment.  There are set-many coordinates, even when $T$ itself is not
well-ordered.  By the monotonicity requirement in
Definition~\ref{calc:natural-layout}, a pointwise increase embeds the old
layout and its designated terms into the new one.  Unequal comparisons
remain unequal because diagrams are preserved by the induced arrows.

For a fixed frozen pattern, apply
Lemma~\ref{lev:ordinal-vectors} to the upward closure, among all ordinal
vectors, of its valid unequal configurations.  It gives a set of
generators, which can be chosen from valid unequal configurations.  Each
generator has a set-sized parameter order.  Take a cardinal above the
sizes of the terms on all these orders, then above the resulting bounds
for all the frozen patterns.  An unequal configuration contains one of
the generating configurations and is consequently witnessed on at most
that many parameters.  Increasing the bound monotonically gives $g$.
This is precisely the ordinal-vector proof of
Proposition~\ref{lev:testing}; its required ingredients here are natural
diagrams, their preservation, and the explicit monotonicity of the
layout, all verified above.  Taking the same supremum over a set of
descriptors proves the last assertion.
\end{proof}

\begin{lemma}[Compact placement for a scheme]
\label{calc:compact-placement}
Let $P$ be any chosen parameter set in a well-ordered presentation, with
one fixed term representation for each parameter.  Let $v$ have length
at most $\sigma$, and let
\[
 \chi>b+\sigma,\qquad \chi^\sigma=\chi.
\]
Suppose that coherent diagrams for $(v,X)$ are prescribed for every
$X\subseteq P$ of size at most $\chi$, and that each individual
prescription has a placement.  At a sufficiently closed eligible level
containing the data, one order enlargement realizes all these
prescriptions simultaneously.  The same assertion holds with bounded
frozen data, after adding its size to $\sigma$.

For a padded order that need not be well-ordered, it suffices instead to
have locally possible coherent prescriptions on every bounded subset
and closure under the profile-count growth used in the proof.  The
conclusion then holds simultaneously at any one specified bounded test
size.  All the bounds are uniform in the fixed scheme and layout data.
\end{lemma}

\begin{proof}
For $X\subseteq P$, let $S_X$ be the union of its chosen supports.  A
profile records the labels of $v$, its formal ordered support, and all
comparisons and coincidences with $S_X$.  If $|X|\le\chi$, there are at
most $\chi$ profiles.  In each strict slot over $S_X$, all the new
positions there form one cluster.  A profile has a continuation over the
whole old order with no new event: place each cluster intact into a
single strict slot of the old order extending its slot over $S_X$.
Order enlargement supplies positions if necessary.  A new event on
enlarging $X$ means a new hit of an old position or a separation of two
previously unseparated new positions by an old position.

Two continuations with no new event over $S_Y$, for $Y\supseteq X$,
have the same diagram over $Y$.  Indeed, their differences consist of
moving whole clusters past blocks of parameter supports.  The regions
belong to distinct slots, so the block transpositions in
Lemma~\ref{calc:parameter-orbits} apply simultaneously.

If simultaneous realization fails, choose $X_0$ of size at most $\chi$.
For each profile over $X_i$, choose a no-new-event continuation and one
prescribed query where it fails.  Adjoin those queries to $X_i$ to obtain
$X_{i+1}$; at limits take unions.  There are at most $\chi$ profiles,
each chosen query has size at most $\chi$, and $\sigma^+\le\chi$.
Thus every $X_i$, including $X_{\sigma^+}$, has size at most $\chi$.
Choose a placement realizing the prescription on $X_{\sigma^+}$.
At every earlier step this placement must have a new event, since
otherwise the no-event comparison would transfer the chosen failure to
it.  There are at most $\sigma$ possible hits and at most $\sigma$
pairs of new positions.  Each event can occur for the first time only
once.  This contradicts the $\sigma^+$ steps.

In a non-well-ordered chart, the number of profiles over $X_i$ is at most
$2^{|X_i|+\sigma+b}$.  Starting with the desired test size, use this as
the next cardinal bound and include a failed query of at most that
original test size for every profile.  The larger sets $X_i$ have
prescriptions because those were supplied on all bounded subsets.
More explicitly, start with an infinite $\beta_0$ above the desired
test size and set $\beta_{i+1}=2^{\beta_i+\sigma+b}$, taking suprema
at limits.  Require the single bound
$\sup_{i\le\sigma^+}\beta_i$ to be strictly below the level.
This is closure under the whole indicated recursion, rather than
merely closure under each individual exponentiation.  It keeps the
final set bounded even though the level has countable cofinality.
Its prescription
is available by hypothesis, and the identical event count is a
contradiction.  This is a set recursion depending only on the starting
arity and fixed data.
\end{proof}

For later reference, fix a convenient explicit threshold for the whole
object layout.  Increase $g$ so that $g(\sigma)\ge b+\sigma$, and put
\begin{equation}
\label{calc:test-threshold}
 c(\sigma,\tau)=2^{g(\sigma)+b+\sigma+\tau},
 \qquad c(\sigma)=c(\sigma,\aleph_0).
\end{equation}
These cardinals exceed the displayed arities and satisfy
$c(\sigma,\tau)^\sigma=c(\sigma,\tau)$.  We close levels under these
functions and their successors, as well as under the diagram and rich
order bounds at the resulting arities.  For a different fixed natural
layout, use its testing function in the same formula.

\begin{lemma}[Placed comparisons of actual tuples]
\label{calc:placed-comparisons}
Suppose $A\to N$ is an arrow between same-level objects and $v$ is a
bounded tuple in $N$.  Its diagrams over bounded subsets of $A$ have a
placed comparison over a well-ordered presentation of $A$.  This
comparison is unique as a diagram over all of $A$.  Two such actual
tuples have equal comparisons if and only if their diagrams agree over
queries in $A$ of size at most $g(|v|+\aleph_0)$.
The conclusion also holds over any natural parameter layout for which
the same coherent local diagrams are given.
\end{lemma}

\begin{proof}
A single bounded request is possible in a presentation: the parameter
diagram is preserved by the arrow, and saturation realizes its desired
extension over the actual parameters in a level object.  Apply
Lemma~\ref{calc:compact-placement} at
$\chi=c(|v|+\aleph_0,\tau)$, for any desired bounded query size $\tau$.
Two outputs agree on all tests of size at most $g(|v|+\aleph_0)$, and
Proposition~\ref{calc:testing} identifies their full placed comparisons.
For any bounded query, choose $\tau$ large enough to include it.  The
unique comparison therefore computes every actual bounded diagram.
This argument does not require an amalgamation relation on types.
\end{proof}

\begin{theorem}[Amalgamation and types for a scheme]
\label{calc:scheme-amalgamation}
At sufficiently closed levels, $\mathcal C$ has amalgamation.  If
$A\to B_0,B_1$ are arrows and $v_i$ are equally indexed bounded tuples
in $B_i$, then an amalgam over $A$ identifying $v_0$ and $v_1$ exists
if and only if their diagrams agree on all parameter queries in $A$ of
size at most $g(|v_i|+\aleph_0)$.

Consequently $\tp(v/A;B)$, defined by identification in an amalgam over
$A$, is an equivalence class.  It is determined by these bounded
diagrams and has an order placement.  Restriction of a type means
restriction of its actual tuple and parameter diagrams.
\end{theorem}

\begin{proof}
Write arrows as inclusions by taking renamed copies.  We prove the
stronger construction with a prescribed identification of bounded
tuples; omit the tuples for ordinary amalgamation.  Let the target
initially be $N_0=B_1$.  Choose increasing bounded lists $d_n$ exhausting
$B_0$, starting with $d_0=v_0$ (or the empty list when no identification
is prescribed).  In each vertex also choose an
increasing countable exhaustion by strong submodels of cardinality less
than the level, and include their enumerations in the $d_n$'s.  Such
exhaustions exist by the component LS property, coherence and
$\cf(\mu)=\omega$.

We construct targets $N_n$ and compatible images of $d_n$.  The invariant
is equality of all bounded joint diagrams of $(d_n,A)$ with their
original diagrams in $B_0$.  For the prescribed initial tuple this is
the hypothesis and Lemma~\ref{calc:placed-comparisons}; the empty initial
tuple is automatic.

Suppose an image of $d_n$ has been chosen.  In a well-ordered chart of
$N_n$, regard $A$ and this whole image as the parameter set $P$.
For a bounded subset of $P$, its required joint diagram with $d_{n+1}$
is read in $B_0$, interpreting the old image as $d_n$.  This prescription
is well-defined on repetitions and coincidences by the invariant.
It is locally possible: equality on the old parameters and saturation
give a realization in $N_n$.  Choose a test cardinal of the form
\eqref{calc:test-threshold}, large enough also for the length of $d_n$,
and apply Lemma~\ref{calc:compact-placement}.  The resulting tuple lies
in an order enlargement $N_{n+1}$ and retains the old assignments.

To verify the invariant, use a natural well-ordered chart of the
\emph{source} $A$.  Both the desired tuple and the newly placed tuple
have comparisons there by Lemma~\ref{calc:placed-comparisons}.  Their
agreement on the chosen tests includes the testing bound for the joint
arity, so Proposition~\ref{calc:testing} makes their comparisons equal.
They therefore agree on every bounded query in $A$.  We have not
applied uniform layout testing to the possibly unnatural image of $A$
in the target chart; compact placement alone was used on that image.

Take the union of the countable target chain.  Its total size is $\mu$,
so the union and upper-bound clauses of the package give an object and
an arrow from $B_1$.  The assignments define a map from $B_0$ into this
union.  They match diagrams on increasing lists that exhaust $B_0$ and
include the component strong-submodel exhaustions.  The image of each
such small submodel is strong: the matching diagram of its enumeration
is implemented by an isomorphism between its original containing
object and the target object.  The arrow criterion
of the package consequently makes this map an arrow, including any
extra basis condition already checked for the current scheme.  The
initial tuples have the same image.  Necessity of diagram equality
follows from preservation in any amalgam.

Amalgamation makes the identification relation transitive: amalgamate
two witnesses over their common intervening object.  Reflexivity and
symmetry follow from identity maps and interchange of the two legs.
The equivalence and testing descriptions now follow from the proved
criterion and Lemma~\ref{calc:placed-comparisons}.
\end{proof}

\begin{corollary}[Realization into a larger level]
\label{calc:mixed-realization}
Let $\nu<\mu$ be sufficiently closed eligible levels.  If $A\to B$ are
$\nu$-objects and $A\to N$ is an arrow into a $\mu$-object, then there
is an arrow $B\to N$ extending the given arrow on $A$.  One may prescribe
the image of a tuple of length less than $\nu$ whose diagrams over
queries in $A$ below $\nu$ already match.  Types of tuples of length
less than $\nu$ over $A$ obtained in a larger object are actual
$\nu$-types, and the mixed diagram comparisons at those arities agree.
\end{corollary}

\begin{proof}
Use the countable exhaustion in the preceding proof, with all source
lists and testing thresholds below $\nu$.  Perform the compact
placement in a chart of $N$.  The placed tuple, all of $A$, and the
previously assigned list together have cardinality at most $\nu$, which
is bounded at $\mu$.  Saturation at $\mu$ therefore realizes their
whole joint diagram back in the actual $N$, fixing the existing
parameters.  Continue there, without enlarging the final target.
For component strongness, first place the source object at level $\mu$
by a presentation arrow.  Mixed preservation and an isomorphism of
$\mu$-objects implementing the matching enumeration show that every
small source strong submodel has strong image in $N$.  The exhaustion
criterion now gives the arrow.

To interpret a tuple of length less than $\nu$ over $A$ inside a larger
object, take a $\nu$-subobject containing $A$ and the tuple, using mixed
LS.  Coherence makes it an extension of $A$.  Mixed preservation
identifies its diagrams with those in the larger object.  The testing
criterion shows independence of the chosen subobject.  This also
justifies the source-level type comparisons used during the construction.
\end{proof}

\subsection{Approximations and stationary free extensions}

Amalgamation has made types over a whole object available.  We now choose
their free extensions by means of sequences in the base: on each bounded
parameter query, almost all entries of the sequence give the same answer.
The next two results construct these sequences and show that their answers
determine an actual type, independently of the choices.

\begin{definition}[Approximation sequence]
An approximation sequence for $p=\tp(v/A)$ at query size $\tau$ is an
indexed family $I=(v_i:i<\kappa)$ in $A$, each $v_i$ having the same
index set as $v$, such that the concatenation of the family is bounded
and, for every $Z\subseteq A$ of size at most $\tau$,
\begin{equation}
\label{calc:exception-estimate}
 \bigl|\{i<\kappa:D(v_i,Z)\ne p\mathord\upharpoonright Z\}\bigr|
 \le |Z|+\aleph_0.
\end{equation}
Here $p\mathord\upharpoonright Z$ denotes its joint diagram, and a
sequence is always indexed: different indices may carry the same tuple.
We use the same terminology for all projections of the displayed tuple.
Whenever $\kappa>|Z|+\aleph_0$, the unique value assumed outside at most
$|Z|+\aleph_0$ indices is called its \emph{majority value} on $Z$.
\end{definition}

\begin{lemma}[Uniform majority approximations]
\label{calc:approximations}
Let $|v|\le\sigma$ and let $\tau\ge\aleph_0$ be bounded.  At
sufficiently closed levels there are approximation sequences for
$p=\tp(v/A)$ at size $\tau$, of every prescribed bounded infinite length
$\kappa$.  They satisfy \eqref{calc:exception-estimate} simultaneously
for all subtuples of $v$.  The construction has bounds depending only
on $b,\sigma,\tau,\kappa$, rather than on the level or the type.

In particular, take $\tau=c(\sigma)$ and
$\kappa=c(\sigma)^+$.  The concatenated sequence then has a uniform
bounded length depending only on $\sigma$ and the fixed scheme.  Its
diagram is a code for $p$ in the precise transport sense proved in
Theorem~\ref{calc:free-extension} below.
\end{lemma}

\begin{proof}
By the rich-order and padding constructions of
Lemmas~\ref{lev:rich-orders} and~\ref{lev:directed-unions}, present $A$
on a padded dense order sufficiently saturated for every support size
used below.  The presentation is an object of the current scheme;
categoricity in the package gives an isomorphism with $A$.
Lemma~\ref{calc:placed-comparisons}, or its compact-placement proof in
the padded chart, places $p$ correctly at size $\tau$.

Keep every extra supporting position that belongs to the old chart.
Partition the other supporting positions into clusters according to
their strict cut over the entire old order.  For every pair of distinct
cuts choose an old position separating them.  Include the shared
supporting positions as well.  This gives a bounded anchor set of size
at most $\sigma+\aleph_0$.  Each unshared cluster now belongs to a
different strict slot over the anchors.

In each such anchor slot, choose $\kappa$ pairwise disjoint convex
regions in the old order and, in the $i$th region, put a copy of the
corresponding cluster with its internal order.  The prescribed
richness supplies these patterns; its required arity is bounded in
terms of $\sigma+\kappa$.  Use the original labels and the unchanged
shared positions to obtain $v_i\in A$.

Choose a query $Z$ and one finite support for each of its elements.
If these supports avoid every region used for the $i$th copies, then
each copied cluster and the original cluster can differ relative to
the query only by passing whole blocks of parameter positions.  There
are no hits or separations inside a cluster.  Their anchor slots are
disjoint, so Lemma~\ref{calc:parameter-orbits} gives
$D(v_i,Z)=p\mathord\upharpoonright Z$.  Every supporting position of
$Z$ belongs to at most one of the regions, and therefore excludes at
most one index $i$.  Finitarity of supports gives at most
$|Z|+\aleph_0$ excluded indices.  The same set of good indices works
for every subtuple of $v$.  Only bounded support patterns, the fixed
labels, rich orders at bounded arities, and the compact-placement
bounds were used.  This proves the uniformity assertions.
\end{proof}

\begin{theorem}[Stationary free extension]
\label{calc:free-extension}
Let $A\to B$ be an arrow at a sufficiently closed level and let
$p=\tp(v/A)$, where $|v|\le\sigma$.  There is a specified actual type
$p^B$ over $B$, called its \emph{free lift}, extending $p$.  At every
desired query size whose approximation and testing costs are bounded at
the originating level, it is computed by transporting sufficiently long
approximations in $A$.  At one level this includes every bounded query
size; in a mixed-level use the costs must remain below the level of $A$.
It has the following properties.
\begin{enumerate}
\item The answer is independent of the approximation sequence, its
length, and the auxiliary test size used in its construction.
\item Free lifts commute with restriction of the tuple, are invariant
under isomorphisms, and satisfy
$(p^B)^C=p^C$ whenever $A\to B\to C$.
\item If $A$ has an eligible size $\nu$ and $B$ has a larger eligible
size, the construction still works for $|v|<\nu$, with its construction
thresholds and approximation lists below $\nu$.  The restriction to
$A$ is the actual $\nu$-type of
Corollary~\ref{calc:mixed-realization}.
\item A sequence from Lemma~\ref{calc:approximations} at the fixed
threshold $c(\sigma)$ is bounded code data.  If an equally indexed
sequence in any sufficiently closed eligible object has the same
diagram, it determines a unique corresponding type there by the same
majority prescription.  Arrows carrying one such sequence to another
carry the corresponding type to its free lift.  In particular an
isomorphism carrying the code sequence transports the type.
\end{enumerate}
We say that an actual type over $B$ is \emph{free from $A$} when it is
the free lift of its restriction to $A$.
\end{theorem}

\begin{proof}
First fix a query bound $\tau$ and take an approximation sequence
$I=(v_i:i<\kappa)$ in $A$ with $\kappa>\tau+\aleph_0$.  Given a
query $Z$ in $B$ of size at most $\tau$, saturation moves it into $A$
over the entire sequence: there is $Z_0\subseteq A$, with the same
enumeration, such that
\begin{equation}
\label{calc:query-move}
 D(I,Z)=D(I,Z_0).
\end{equation}
The total length $|I|+|Z|$ is bounded.  For mixed levels it is bounded
at the originating level, and the mixed saturation clause applies.
Define the answer on $Z$ to be the majority of $D(v_i,Z)$, equivalently
the answer $p\mathord\upharpoonright Z_0$.  Equation
\eqref{calc:query-move} preserves every entry of the sequence, so this
definition is independent of the chosen $Z_0$ and retains the same
exception estimate.

If $I'$ is another sufficiently long approximation sequence, move $Z$
into $A$ over $I\cup I'$.  On that moved query both majorities compute
the original $p$.  Thus the answers agree.  This proof requires
saturation over the concatenation of the sequences; it does not
require either approximation estimate to hold for a query containing
that concatenation.  The averaging queries still have size at most
$\tau$.  The identical argument compares different query bounds on
their common domain.  It also proves coherence under restrictions and
reindexing of queries.

We next obtain an actual type.  Use
$\tau=c(\sigma)$ and $\kappa=\tau^+$.  In a well-ordered chart of $B$
prescribe the majority answers on every query of size at most $\tau$.
For each single query these are realized by some $v_i\in A\subseteq B$,
since fewer than $\kappa$ indices are exceptional.  Thus they have
local placements.  Lemma~\ref{calc:compact-placement} gives a placement
realizing all of them in an order enlargement of $B$.  This is an
actual tuple in an extension object, and
Theorem~\ref{calc:scheme-amalgamation} makes its type unique, since
$\tau\ge g(\sigma)$.  This type is $p^B$.  Its restriction to $A$
equals $p$ by the same testing criterion.

For a larger desired bounded query size $\tau'$, repeat with a power
closed bound at least $c(\sigma,\tau')$.  The comparison of sequences
just proved identifies the resulting type with $p^B$ on the initial
testing threshold and hence everywhere.  In a mixed-level use, this
additional direct majority evaluation is asserted only when its costs
remain below the originating level.  The type itself is already
actual over $B$, and its diagrams at larger arities are those of this
actual type; no averaging inside a smaller model at unbounded arities
is being asserted.

Restriction to a subtuple commutes with every majority evaluation.
Choose the threshold large enough for both arities and use testing to
obtain equality of types.  Isomorphism invariance follows by applying
the isomorphism to the sequences and their query diagrams.
For transitivity, let $J\subseteq B$ be a canonical approximation
sequence for $p^B$ at a sufficient fixed threshold, and retain an
original sequence $I\subseteq A$ at that threshold.  The latter
computes $p^B$ on queries in $B$ by the construction.  Move a query in
$C$ into $B$ over $I\cup J$.  The two sequences give the same answer
there, because each computes $p^B$.  The first also computes $p^C$,
and the second computes $(p^B)^C$.  Testing proves transitivity.
All these sequence lengths can be chosen as the uniform bounds in
Lemma~\ref{calc:approximations}, below the smallest originating level.
This proves the mixed-level assertions as well.

Finally, suppose $J\subseteq B$ has the same whole diagram as a fixed
code $I\subseteq A$ at threshold $c(\sigma)$.  For each query $Z$ in
$B$ of that size, saturation realizes $D(J,Z)$ over $I$ in $A$.
The resulting majority is therefore well-defined, with the same
exception estimate.  For each query one of the entries of $J$ realizes
the answer in $B$.  Coherence, compact placement and testing give an
actual unique type, as in the preceding construction.  Moving queries
over code sequences shows that an arrow carrying $J$ into a larger
object gives exactly its free lift there.  The code length and every
threshold used here depend only on $\sigma$ and the fixed scheme,
which proves the final assertion.
\end{proof}

\begin{lemma}[An average criterion for descent]
\label{calc:average-descent}
Let $A\to B$ have eligible sizes $\nu\le\mu$, and let $q=\tp(d/B)$
be actual, with $|d|=\sigma<\nu$.  Choose
$\chi\ge c(\sigma)$ below $\nu$.  Suppose that a bounded sequence
$I=(d_i:i<\kappa)$ in $A$, with $\chi+\aleph_0<\kappa<\nu$, has
majority value $q\mathord\upharpoonright Z$ on every query
$Z\subseteq B$ of size at most $\chi$, with at most
$|Z|+\aleph_0$ exceptions.  Then $q$ is the free lift of its actual
$\nu$-restriction to $A$.
\end{lemma}

\begin{proof}
Let $p=q\mathord\upharpoonright A$ and choose a canonical sequence
$J\subseteq A$ approximating $p$ at size $\chi$.  For a query
$Z\subseteq B$ of that size, mixed saturation moves it to $Z_0$ in
$A$ over the whole $I\cup J$.  The total length is less than $\nu$.
The majority from $I$ on $Z$ computes $q$, and on $Z_0$ computes $p$.
These answers agree by diagram preservation.  The majority from $J$
on $Z_0$ also computes $p$, and on $Z$ computes $p^B$.
Hence $q$ and $p^B$ agree on all $\chi$-queries.  Apply the testing
criterion.  The same proof works for an arbitrary longer indexed
family whenever a bounded subsequence of the displayed length retains
the exception estimate.
\end{proof}

\subsection{Products and local character}

\begin{proposition}[Products of bounded types]
\label{calc:products}
For bounded types $p(x),q(y)$ over $A$, realize $p$ in an extension
$B\ge A$ and then realize $q^B$ in a further extension.  The joint
type of $(x,y)$ over $A$, denoted $p\otimes q$, is independent of all
choices.  Binary products are symmetric, with the indicated
permutation of coordinates.  Finite products associate, are
permutation invariant, and commute with coordinate projections.
Furthermore
\begin{equation}
\label{calc:free-product-formula}
 (p\otimes q)^B=p^B\otimes q^B.
\end{equation}
These assertions hold across eligible levels when all originating
arities and uniform testing costs are bounded at the smallest level.

If $x/B$ is free from $A$ and $x,d$ have product type over $B$, then
$x$ and $(d,e)$ have product type over $A$ for every bounded tuple
$e$ from $B$.
\end{proposition}

\begin{proof}
Two realizations of $p$ can be identified over $A$ by amalgamation.
In a common model containing their witnessing extensions, the free
lifts of $q$ agree by stationarity and restriction.  A second
type-identifying amalgamation identifies the realizations of $q$.
This proves independence of the choices in the definition.

We give the cardinal calculation for symmetry.  Fix a bounded query
$Z$ in $A$, and put
$e=|Z|+|x|+|y|+\aleph_0$.  Choose approximation sequences
$(a_i:i<\kappa)$ for $p$ and $(b_j:j<\lambda)$ for $q$, sufficient
for queries of size $e$, where
\[
 \kappa=e^+,\qquad \lambda=\kappa^+.
\]
All these cardinals lie below a sufficiently closed level.
For the array $M_{ij}=D(a_i,b_j,Z)$, first realize $x$ and freely
realize $y$.  The resulting answer $R$ is obtained by first taking
the majority in $i$ and then in $j$.  Except at at most $e$ columns,
the column majority is $R$, and in each such column at most $e$
entries differ from $R$.  Reversing the realization order gives an
answer $S$: except at at most $e$ rows, its row majority is $S$,
with at most $e$ exceptions in each such row.  The bounds concern
the size of $(x,Z)$ or $(y,Z)$, not the lengths of the sequences.

There are $\kappa$ good rows.  Their exceptional columns together
number at most $\kappa e=\kappa$.  Since $\lambda>\kappa$, choose
a column outside this set and outside the at most $e$ bad columns.
Every good-row entry in that column equals $S$, whereas all but at
most $e$ entries equal $R$.  As $\kappa>e$, some entry has both
values.  Thus $R=S$.  This explains why two equally sized sequences
would not suffice for the argument.  Testing gives symmetry of
the joint types.

For \eqref{calc:free-product-formula}, take the two sequences just
used in $A$ and a sequence $K\subseteq A$ approximating the actual
old product.  Choose their accuracy at a common sufficient test
threshold.  Given a query $Z$ in the enlarged base $B$, move it into
$A$ over these three whole sequences.  Both successive free
evaluations over $B$ are the iterated majority of the same array
$D(a_i,b_j,Z)$.  Transitivity permits the original sequence for the
second factor to be used after the first factor has been realized.
The query movement identifies this array with the old array in $A$,
where it computes $p\otimes q$.  The sequence $K$ evaluates that
old product and computes its free lift on the original query.
The answers therefore agree.  Testing proves the formula.  All the
lists fixed during the query movement are bounded, including in the
mixed case.

Now realize $x,y,z$ successively along a chain of free lifts from
$A$.  The first two form $p\otimes q$, and $z$ is free over a model
containing them, so the joint type is $(p\otimes q)\otimes r$.
Over a model containing $x$, the last two have type
$q^B\otimes r^B=(q\otimes r)^B$, by the formula just proved.
Thus the same joint type is $p\otimes(q\otimes r)$.  This proves
associativity.  Projection follows from the projection property of
free lifts in each successive realization.  Associativity and the
binary symmetry permit every permutation of a finite list of factors.

For the final assertion, use symmetry and a type-identifying amalgam
to realize the given joint type by first putting $d$ in a model
containing $B$ and then realizing $x$ freely over it from $B$.
Since $x/B$ is itself free from $A$, transitivity makes this last
realization free from $A$.  The tuple $(d,e)$ is in the intervening
model, so the definition of product gives the assertion.  The amalgam
identifies the full original tuple over $B$, and hence retains its
actual joint type over $A$.
\end{proof}

\begin{proposition}[Local character and continuity]
\label{calc:local-character}
Let $(A_i:i<\delta)$ be an increasing chain of objects at one eligible
level $\mu$, whose union $A$ has size $\mu$, and let $\delta$ be a
nonzero limit ordinal.  If $p=\tp(v/A)$ and either $v$ is finite or
$|v|<\cf(\delta)$, then $p$ is free from some $A_i$.
If, in addition, every restriction $\tp(v/A_i)$ is free from a fixed
base $X\leq A_0$, then $p$ is free from $X$.
\end{proposition}

\begin{proof}
Pass to a cofinal chain indexed by the regular cardinal
$\gamma=\cf(\delta)$.  Choose a sufficient testing size
$\chi\ge c(|v|+\aleph_0)$ below $\mu$.  If $\gamma$ is larger than
the size of an approximation sequence for $p$ at that threshold,
the whole sequence lies in one stage.  Its exception estimate is
still correct on all queries in $A$, so
Lemma~\ref{calc:average-descent}, at the same level, applies.

In the other case $\gamma<\mu$.  Choose a regular
$\kappa<\mu$ above $\gamma,\chi,|v|$ and the needed bounds, and
take $\kappa$ approximations inside $A$.  Every whole approximation
tuple belongs to some stage: this uses finiteness, or the hypothesis
$|v|<\gamma$.  If every stage contained fewer than $\kappa$ of the
tuples, regularity of $\kappa$ and $\gamma<\kappa$ would make their
union contain fewer than $\kappa$ tuples.  Thus one stage contains
$\kappa$ of them.  Pass to those indices.  The exception estimate
survives unchanged, and Lemma~\ref{calc:average-descent} again
identifies $p$ as the free lift from that stage.

For the final assertion, local character makes $p$ free from some
$A_i$.  Its restriction there is free from $X$ by hypothesis.
Transitivity in Theorem~\ref{calc:free-extension} makes $p$ free
from $X$.
\end{proof}

\subsection{Universal strips and approximate retractions}

The remaining construction provides maps from an extension back into its
base which fix each individual base element at all but finitely many
indices.  It
will allow isolation to descend from a larger level without requiring a
bounded tuple to lie in one stage of a countable exhaustion.

\begin{lemma}[Closure of a bounded support set]
\label{calc:support-closure}
Let $J\subseteq J^*$ be linear orders, let $V$ be a convex strip in
$J^*$, and put $J_0=J\setminus V$.  Allow any fixed finite coloring
of positions in $J_0$, including their side of $V$.  Let
$T^*\subseteq J^*$ and $S_0\subseteq J_0$ be bounded sets.  Given
an infinite $\chi$, choose an infinite $\lambda$ with
\[
 \lambda\ge\chi+|T^*|+|S_0|,
 \qquad \kappa=2^\lambda.
\]
There is $S\subseteq J_0$, of size at most $\kappa$, containing
$S_0\cup(T^*\cap J_0)$, with this property: for every
$R\subseteq S$ of size at most $\chi$ and every indexed tuple $Q$
of at most $\chi$ positions from $J_0$, there is an equally indexed
tuple $Q'$ in $S$ having the same order, equality and color pattern
over $R\cup T^*$.  In particular all positions named by $R$ are
fixed whenever they occur in $Q$.

At a strong limit level with the displayed input bounded, $S$ is
bounded.  The assertion needs no well-order hypothesis on $J$.
\end{lemma}

\begin{proof}
Build $S_\alpha$ for $\alpha<\chi^+$, starting with the required
set.  At a successor stage, for every $R\subseteq S_\alpha$ of
size at most $\chi$, and every order, equality and color pattern of
at most $\chi$ further positions over $R\cup T^*$ that is actually
realized in $J_0$, choose one realizing tuple and adjoin its positions.
Use unions at limits.  If $|S_\alpha|\le\kappa$, the number of
choices of $R$ is at most $\kappa^\chi=\kappa$, and the number of
patterns for each $R$ is at most $2^\lambda=\kappa$.  Each chosen
tuple adds at most $\chi$ positions.  Consequently every stage and
their union have size at most $\kappa$.

For $R$ contained in the final union, regularity of $\chi^+$ puts
all of $R$ into a single stage.  The actual tuple $Q$ witnesses that
its pattern was among the requirements processed at the next stage.
Its chosen witness is the required $Q'$.  Equality with elements of
$R$ and $T^*$ is part of the pattern, so it is preserved literally.
This also covers arbitrary enumerations and repetitions in $Q$.
\end{proof}

\begin{proposition}[Universality of a fresh strip]
\label{calc:universal-strip}
Let $\mathcal C$ already satisfy the working package.  In addition to
its fixed template $F$, one may use a presentation $G$ with the
following properties: it has a fixed set of finite-support labels;
every presentation $G(J)$ used below, for an order of size $\mu$, is
an object of $\mathcal C$ at level $\mu$; order embeddings give
$\mathcal C$-arrows; and the diagrams of its term lists are natural,
including in the order enlargements and padded orders used.  Assume
also that these labels represent every element of $G(J)$.
Require $\mu$ to exceed the cardinality of the label set of $G$ and
to be closed under the bounds below with that cardinal included in
their fixed input data.

Put $A=G(J)$.  Insert into $J$ a fresh convex strip of size $\mu$
with countably many disjoint convex substrips, each of size $\mu$,
and denote the resulting object by $A^+$.  Then every same-level
extension of $A$ embeds into $A^+$ over $A$.  In particular this
applies to $G=F$ and to well-ordered presentations.  No separate
union or categoricity assertion for a category of $G$-presentations
is required.
\end{proposition}

\begin{proof}
We first realize a bounded type $p(v/A)$ using just one fresh
substrip $V$.  Choose a testing threshold $\chi\ge c(|v|+\aleph_0)$
large enough also for $|v|$ and the number of labels of $G$.
Take a sequence $I$ in $A$ approximating $p$ at query size $\chi$.
Represent its entries by $G$-terms and collect their old supports in
$S_0\subseteq J$.  Apply Lemma~\ref{calc:support-closure}, with
the side of $V$ recorded, to obtain a bounded $S\subseteq J$.
Its closure permits every at-most-$\chi$ old query pattern to move
into $S$ over any at-most-$\chi$ subset of $S$.  The approximation
supports, which may be larger than $\chi$, can instead be included
as the fixed set $T^*$ in that Lemma.  Thus all of them are fixed
during every such movement.

Let $P_S$ be the set of all $G$-terms on $S$.  It is bounded, since
$G$ has a bounded set of finite-support labels.  The natural
parameter diagram of $P_S$ is the same in $A$ and in $G(S\cup V)$.
The latter is a level object because $|V|=\mu$.  Saturation there
realizes the desired whole joint diagram of $(v,P_S)$, giving a
tuple $w$ whose supporting positions belong to $S\cup V$.
Only at most $|v|+\aleph_0\le\chi$ of its supporting positions
are old; call them $R\subseteq S$.

For any query $Z$ in $A$ of size at most $\chi$, move its chosen
support tuple into $S$ over $R$ and the fixed approximation
supports, preserving the side of $V$.  Let $Z'$ be the resulting
term tuple.  All comparisons with the new supporting positions
of $w$ in $V$ are forced by that side.  Naturalness therefore gives
\[
 D(w,Z)=D(w,Z'),\qquad D(I,Z)=D(I,Z').
\]
The majority from $I$ computes $p$ on both queries.  The chosen
realization of $(v,P_S)$ gives $D(w,Z')=p\mathord\upharpoonright Z'$.
It follows that $D(w,Z)=p\mathord\upharpoonright Z$.
Testing of actual types, using a natural $F$-chart of $A$ if $G$
differs from $F$, proves $\tp(w/A)=p$.  Naturality of $G$ was used
only for the displayed support movements; the testing and saturation
are those already proved in $\mathcal C$.

Now let $B\ge A$ be arbitrary and choose increasing bounded lists
$d_n$ exhausting it, including countable strong-submodel
exhaustions in every vertex.  Use the first substrip to realize the
type of $d_0$ over $A$.  At subsequent step $n$, the target base consists
of $A$ and the whole previously used substrips.  Suppose $d_n$ has
already been assigned there with its correct type over $A$.
Theorem~\ref{calc:scheme-amalgamation} amalgamates $B$ with this
target base over $A$, identifying that entire assigned tuple.
In this amalgam, $d_{n+1}$ has an actual bounded type over the
target base.  Apply the one-substrip construction to the next fresh
substrip.  Since its type includes equality with the already
assigned coordinates, the new assignment extends the previous one.
The final union of the assignments lands in $A^+$.  The arrow
criterion applies to the included strong-submodel exhaustions, so
it is an arrow $B\to A^+$ over $A$.
\end{proof}

\begin{corollary}[Approximate retractions]
\label{calc:retractions}
For same-level $A\le B$ there are arrows
$(r_i:B\to A:i<\mu)$ such that, for every $a\in A$,
\[
 \bigl|\{i<\mu:r_i(a)\ne a\}\bigr|<\aleph_0.
\]
The maps are not asserted to fix all of $A$ at once.
\end{corollary}

\begin{proof}
Present $A$ by $F$ on a well-order consisting of $\mu$ disjoint
convex blocks $J_i$, each of cardinality $\mu$ and itself containing
countably many convex substrips of that size.  The whole order still
has cardinality $\mu$, and the working package identifies its
presentation with $A$.  Let $A_i$ be the presentation on the
complement of $J_i$.  It is a level subobject of $A$.
Proposition~\ref{calc:universal-strip}, with $J_i$ as the inserted
strip, embeds $B$ into $A$ over $A_i$; call the arrow $r_i$.
Choose one finite supporting set for $a$.  It meets only finitely
many $J_i$.  For every other index, $a\in A_i$ and $r_i(a)=a$.
\end{proof}

\begin{lemma}[Descent through approximate retractions]
\label{calc:retraction-descent}
Let $A\le U$ be objects at an eligible level $\mu$, and let
$A'\le U'$ be objects at an eligible level $\nu\le\mu$, with
compatible arrows $A'\le A$ and $U'\le U$.  Suppose there are
arrows $(r_i:U\to A:i<\nu)$ such that every $a\in A$ is fixed
for all but finitely many $i$, and $r_i[U']\subseteq A'$ for all
$i<\nu$.  Then for every tuple $d$ in $U'$ of length less than
$\nu$, its actual type over $A$ is the free lift of
$\tp(d/A';U')$.  All comparisons are comparisons of the working
diagrams, including mixed-level diagrams at arities bounded at $\nu$.
\end{lemma}

\begin{proof}
Coherence makes $r_i\mathord\upharpoonright U':U'\to A'$ an
arrow: its image and $A'$ are subobjects of $A$, and its image is
contained in $A'$.  Put $\sigma=|d|+\aleph_0$, choose
$\chi\ge c(\sigma)$ below $\nu$, and choose
$\chi^+\le\kappa<\nu$.  Retain any $\kappa$ indices and form
$I=(r_i(d))$ on them.  This is a bounded concatenated list in $A'$.

For a query $Z\subseteq A$ of size at most $\chi$, at most
$|Z|+\aleph_0$ indices fail to fix every element of $Z$.  At every
other index, arrow preservation gives
\[
 D(r_i(d),Z)=D(r_i(d),r_i(Z))=D(d,Z).
\]
Thus $I$ has the required majority estimate for the actual type
over $A$.  On queries in $A'$ the same identity is a comparison at
level $\nu$, using the restricted arrows.  Mixed preservation
identifies it with the larger-level identity.

Apply Lemma~\ref{calc:average-descent} to $A'\le A$.  Its
restriction type is precisely $\tp(d/A';U')$, by
Corollary~\ref{calc:mixed-realization}.  The conclusion is therefore
the stated free lifting of the actual type, including tuples that
span a chosen resolution of $U'$.
\end{proof}

The preceding calculus will be applied only after the working package
has been established for the relevant scheme.  Its approximation,
testing and support-closure bounds are functions of fixed scheme data
and bounded input arities.  This is the uniformity needed when the
schemes are iterated at finite depths.
